% !TeX root = ../../Book3.tex
\section{全纯覆盖}
\label{b3:sec:1102}

全纯覆盖把上节的拓扑条件与复坐标结合。它处处局部双全纯；与存在临界点的分歧映射相比，解析函数的提升因此没有局部角度折叠。

\subsection{全纯覆盖}
\label{b3:subsec:110201}

\term[riemann-qumian]{Riemann 曲面}[Riemann surface]是 Hausdorff、第二可数且具有复平面局部坐标的连通曲面，各坐标变换全纯且导数非零。曲面间映射全纯，指其坐标表达全纯。若一个拓扑覆盖投影到 Riemann 曲面，则在每一层拉回底坐标，给覆盖空间唯一的 Riemann 曲面结构，使投影局部双全纯。坐标交叠的变换直接是底空间的坐标变换，所以相容性完整成立。

\ProseParagraphBreak
在这个结构下甲板变换自动全纯：在一层坐标中它把相同底坐标传到另一层，坐标表达为恒等。全纯函数 \(f\) 的提升 \(f\circ p\) 也全纯。反向，一个甲板不变全纯函数可按底点的任意提升定义底函数，各层逆保证它局部全纯；这就是下降。

\subsection{指数映射}
\label{b3:subsec:110202}

\begin{proposition}\label{b3:prop:11-exponential-cover}
指数 \(p(w)=e^w\) 是 \(\mathbb C\rightarrow\mathbb C^*\) 的万有覆盖，其甲板群为 \(w\mapsto w+2\pi\mathrm{i}n\)、\(n\in\mathbb Z\)。
\end{proposition}
\begin{proof}
在任意目标点的小无零圆盘上取全纯对数 \(L\)。全部原像层为 \(L(V)+2\pi\mathrm{i}n\)，互不交，指数在每层有逆 \(L+2\pi\mathrm{i}n\)，所以是覆盖。源平面单连通，故为万有覆盖。甲板变换满足 \(e^{T(w)-w}=1\)，连续性使 \(T(w)-w\) 是固定的 \(2\pi\mathrm{i}n\)。反向这些平移显然是甲板变换。
\end{proof}
对数在覆盖空间上就是坐标 \(w\)。它不下降为底空间单值函数，因为甲板平移改变它；\(e^{mw}\) 在 \(m\in\mathbb Z\) 时不变，下降为 \(z^m\)。这种下降判别把分支周期转成群不变性。

\subsection{去心区域的覆盖}
\label{b3:subsec:110203}

圆环 \(A(r,R)\) 的指数覆盖限制在 \(\log r<\operatorname{Re}w<\log R\)。去心盘的覆盖限制在左半平面。它们的源域都单连通，故是各自万有覆盖。把左半平面用分式变换映成圆盘，可得到圆盘到去心盘的显式覆盖，但这个映射不是单射；其甲板变换在圆盘坐标中变为抛物型自同构。

\ProseParagraphBreak
有限覆盖 \(w\mapsto w^m:\mathbb C^*\rightarrow\mathbb C^*\) 的甲板群为 \(w\mapsto\zeta w\)、\(\zeta^m=1\)。单位正圈提升只转过 \(2\pi/m\)，走 \(m\) 次才闭合。它的源仍不单连通，因此不是万有覆盖；有限层与万有性是不同条件。

\begin{proposition}[圆环的有限层覆盖]\label{b3:prop:11-annulus-power-cover}
\(z\mapsto z^m\) 把 \(A(r,R)\) 作为 \(m\) 层全纯覆盖映到 \(A(r^m,R^m)\)，目标模量是源模量的 \(m\) 倍。
\end{proposition}
\begin{proof}
目标每一点有 \(m\) 个根，模均落在 \((r,R)\)。导数 \(mz^{m-1}\) 在圆环非零；选足够小的目标圆盘可取 \(m\) 个不同根分支，且它们的像互不相交、覆盖全部原像，故满足均匀覆盖定义。模量由对数半径比计算，即 \(\log(R^m/r^m)/(2\pi)=m\log(R/r)/(2\pi)\)。
\end{proof}
有限层覆盖不保环域模量，因为它并非一一对应的共形等价。第八章模量不变性针对双全纯映射，不能把“全纯且导数非零”自动升级成那个假设。通过条带 \(\log r<\operatorname{Re}w<\log R\) 的指数覆盖看，同一事实是甲板平移周期由 \(2\pi\mathrm{i}\) 的格点变为一个有限指数子群。

\ProseParagraphBreak
对于去心盘，指数覆盖使用左半平面，去心平面则使用整个平面。两者都只有一个可见的去心点，却给不同万有覆盖类型：前者的外边界提供额外限制，后者没有外边界。把所有“去心域”只画成一个孔，容易忽略这个全局差别。

\subsection{与分歧覆盖的区别}
\label{b3:subsec:110204}

映射 \(z\mapsto z^m\) 在含0的平面上有局部次数 \(m\)，0处导数为零，不能作为覆盖。删去分歧点及其值后，才成为前一小节的有限覆盖。

\begin{proposition}\label{b3:prop:11-proper-unbranched}
若非恒常全纯映射 \(f:X\rightarrow Y\) 适当且无临界点，即紧集原像紧，则它是有限层全纯覆盖；若 \(Y\) 连通，层数固定。
\end{proposition}
\begin{proof}
无临界点给局部双全纯，故像开。紧集原像紧使像闭：若 \(f(x_n)\rightarrow y\)，这些像点连同极限构成紧集，原像紧使 \(x_n\) 有子列趋于 \(x\)，故 \(f(x)=y\)。像非空，所以满射。每个纤维既离散又紧，因此有限。围绕有限原像点取互不交逆函数邻域；缩小底邻域后不可能还有额外原像，否则一列额外点像趋于 \(y\)，紧集原像紧给极限落入某个既选邻域，矛盾。故底邻域均匀覆盖，有限层数由连通性固定。
\end{proof}
满足紧集原像紧的映射称为\term[shidang-yingshe]{适当映射}[proper map]；这里的条件是一项全局紧性要求。多项式去掉临界值及其全部原像后满足这些条件，给出典型有限覆盖。

\ProseParagraphBreak
多项式 \(P\) 的适当性可直接从最高次项证明。\(|z|\) 足够大时 \(|P(z)|\ge c|z|^d\)，所以紧目标集的原像有界；连续性又使它闭，故紧。删去有限临界值及其全部原像，\(P\) 就成为 \(d\) 层无分歧覆盖。以 \(P(z)=z^3-3z\) 为例，临界值为 \(\pm2\)，删掉它们后每个目标点有三个不同根；小目标盘内全部根分支由局部逆取得，适当性排除第四种逃向无穷的局部层。重新补回临界值后，局部 \(z^2\) 形状在相应纤维出现，完整映射是分歧覆盖。

\ProseParagraphBreak
另一方面，\(e^z\) 不适当：值1的原像 \(2\pi\mathrm{i}\mathbb Z\) 无界。它仍是覆盖，因为指数的周期让全部无穷层有同一个均匀邻域。适当性是取得有限层覆盖的充分工具，不是覆盖的必要条件。应分别检查导数、适当性、满射与完整层分解，不能用其中一项替代其余。
